
\documentclass{article}
\usepackage{colortbl}
\usepackage{makecell}
\usepackage{multirow}
\usepackage{supertabular}

\begin{document}

\newcounter{utterance}

\centering \large Interaction Transcript for game `cladder', experiment `full\_v1.5\_default', episode 7123 with refineonly{-}Qwen3.5{-}27B.
\vspace{24pt}

{ \footnotesize  \setcounter{utterance}{1}
\setlength{\tabcolsep}{0pt}
\begin{supertabular}{c@{$\;$}|p{.15\linewidth}@{}p{.15\linewidth}p{.15\linewidth}p{.15\linewidth}p{.15\linewidth}p{.15\linewidth}}
    \# & \multicolumn{2}{c}{Player} && \multicolumn{2}{c}{Game Master} \\
    \hline

    \theutterance \stepcounter{utterance}  
    & & & \multicolumn{4}{p{0.6\linewidth}}{
        \cellcolor[rgb]{0.9,0.9,0.9}{
            \makecell[{{p{\linewidth}}}]{
                \texttt{\tiny{[P1$\langle$GM]}}
                \texttt{You are an expert in causal inference. The following question is not a typical commonsense query, but rather a meticulously designed question created by a professor specializing in causal inference, intended to assess the students' mastery of the course content.} \\
\texttt{Imagine a self{-}contained, hypothetical world with only the following conditions, and without any unmentioned factors or causal relationships: Jyka has a direct effect on yupt and kwox. Yupt has a direct effect on kwox. Method 1: We look at how jyka correlates with kwox case by case according to yupt. Method 2: We look directly at how jyka correlates with kwox in general. To understand how jyka affects kwox, is it more correct to use the Method 1 than Method 2?} \\
\texttt{Start your answer with 'Yes' or 'No', followed by additional reasoning or evidence to support your explanation.} \\
            }
        }
    }
     \\ \\

    \theutterance \stepcounter{utterance}  
    & \multicolumn{4}{p{0.6\linewidth}}{
        \cellcolor[rgb]{0.561,0.737,0.561}{
            \makecell[{{p{\linewidth}}}]{
                \texttt{\tiny{[P1$\rangle$GM]}}
                \texttt{No} \\
\\ 
\texttt{In the described causal structure, **Jyka** influences **Kwox** via two pathways: a direct path (\$Jyka \textbackslash{}rightarrow Kwox\$) and an indirect path mediated by **Yupt** (\$Jyka \textbackslash{}rightarrow Yupt \textbackslash{}rightarrow Kwox\$).} \\
\\ 
\texttt{**Method 2** (examining the general correlation between Jyka and Kwox) estimates the **Total Effect**, which is the sum of the direct and indirect effects. Since the objective is to understand "how Jyka affects Kwox" in a comprehensive sense, capturing this total impact is the standard and correct approach.} \\
\\ 
\texttt{**Method 1** (stratifying or conditioning on Yupt) blocks the indirect path (\$Jyka \textbackslash{}rightarrow Yupt \textbackslash{}rightarrow Kwox\$). By conditioning on a mediator, one isolates only the **Direct Effect**. While this is useful for specific mechanistic questions, it fails to capture the full causal influence of Jyka on Kwox because it artificially removes the effect transmitted through Yupt.} \\
\\ 
\texttt{Therefore, Method 1 is not "more correct" than Method 2 for understanding the overall effect; using Method 1 would lead to an underestimation of the total causal impact if the indirect path is non{-}zero. Method 2 provides the complete picture of the causal relationship.} \\
            }
        }
    }
    & & \\ \\

    \theutterance \stepcounter{utterance}  
    & & & \multicolumn{2}{p{0.3\linewidth}}{
        \cellcolor[rgb]{0.9,0.9,0.9}{
            \makecell[{{p{\linewidth}}}]{
                \texttt{\tiny{[GM$|$GM]}}
                \texttt{No} \\
            }
        }
    }
    & & \\ \\

    \theutterance \stepcounter{utterance}  
    & & & \multicolumn{2}{p{0.3\linewidth}}{
        \cellcolor[rgb]{0.9,0.9,0.9}{
            \makecell[{{p{\linewidth}}}]{
                \texttt{\tiny{[GM$|$GM]}}
                \texttt{no} \\
            }
        }
    }
    & & \\ \\

    \theutterance \stepcounter{utterance}  
    & & & \multicolumn{2}{p{0.3\linewidth}}{
        \cellcolor[rgb]{0.9,0.9,0.9}{
            \makecell[{{p{\linewidth}}}]{
                \texttt{\tiny{[GM$|$GM]}}
                \texttt{game\_result = WIN} \\
            }
        }
    }
    & & \\ \\

\end{supertabular}
}

\end{document}
